\chapter{Grundlagen der Optimierung}
\label{kap:grundlagen}

Das in dieser Arbeit betrachtete Problem -- die Zuordnung von Arbeitsgängen zu
Mitarbeitern bei minimaler Gesamtdauer -- ist kein Einzelfall, sondern eine
Ausprägung einer seit Jahrzehnten untersuchten Problemklasse. Dieses Kapitel
ordnet die Aufgabenstellung in die Theorie der Ablaufplanung ein, bestimmt ihre
Komplexität und stellt die drei Lösungswege vor, die im weiteren Verlauf
verglichen werden: ein exaktes Verfahren, zwei Konstruktionsheuristiken und ein
stochastisches Verbesserungsverfahren. Die Darstellung folgt im Wesentlichen
\textcite{pinedo2016scheduling}, der die Ablaufplanung als eigenständige
Disziplin umfassend aufbereitet.

\section{Zuordnungs- und Reihenfolgeprobleme}
\label{sec:einordnung}

\subsection{Ablaufplanung als Problemklasse}

Die \Index{Ablaufplanung} beschäftigt sich mit der Frage, wie eine Menge von
Aufgaben auf eine Menge von Ressourcen verteilt und zeitlich angeordnet wird,
sodass ein vorgegebenes Gütemaß möglichst gut erfüllt ist. Die Ressourcen
heißen in der Literatur traditionell \Index{Maschinen}, die Aufgaben
\Index{Jobs} -- eine Sprechweise, die aus der Fertigungssteuerung stammt und
auch dann beibehalten wird, wenn es sich, wie im vorliegenden Fall, um Menschen
handelt. Diese terminologische Übertragung ist unproblematisch, solange man sich
vergegenwärtigt, dass das Modell von allen Eigenschaften absieht, die über die
Bearbeitungsdauer und die Zulässigkeit einer Zuordnung hinausgehen.

Zu unterscheiden sind zwei Fragestellungen, die in der Praxis oft verschmelzen.
Ein \Index{Zuordnungsproblem} fragt allein danach, \emph{wer} welche Aufgabe
übernimmt; ein \Index{Reihenfolgeproblem} fragt zusätzlich, \emph{wann} das
geschieht. Sind die Aufgaben voneinander unabhängig, existieren keine
Rüstzeiten und keine Fälligkeitstermine, so zerfällt das Reihenfolgeproblem: Die
Reihenfolge, in der ein Mitarbeiter seine zugeteilten Arbeitsgänge abarbeitet,
ist für die Gesamtdauer ohne Belang, weil sich die Summe seiner
Bearbeitungszeiten dadurch nicht ändert. Es bleibt ein reines
Zuordnungsproblem. Genau diese Situation liegt der vorliegenden Arbeit zugrunde,
und sie ist der Grund dafür, dass sich das Problem so kompakt formulieren lässt.

\subsection{Die Drei-Feld-Notation}

Für die Klassifikation von Problemen der Ablaufplanung hat sich die von
\textcite{graham1969bounds} eingeführte und später erweiterte
\Index{Drei-Feld-Notation} durchgesetzt. Ein Problem wird durch ein Tripel
\begin{equation}
  \alpha \mid \beta \mid \gamma
\end{equation}
beschrieben. Das Feld $\alpha$ benennt die \Index{Maschinenumgebung}, das Feld
$\beta$ sammelt Nebenbedingungen und Sonderregeln, und $\gamma$ gibt die
\Index{Zielfunktion} an. Die Notation ist deshalb so nützlich, weil sie eine
Aufgabenstellung in wenigen Zeichen so präzise festlegt, dass sich unmittelbar
in der Literatur nachschlagen lässt, was über sie bekannt ist -- insbesondere,
ob sie effizient lösbar ist.

Für die Maschinenumgebung paralleler Maschinen sind drei Stufen zu
unterscheiden, die in Abbildung~\ref{abb:maschinenmodell} einander
gegenübergestellt sind. Bei \Index{identischen parallelen Maschinen} ($P$)
benötigt jede Maschine für denselben Job dieselbe Zeit $p_j$. Bei
\Index{uniformen parallelen Maschinen} ($Q$) besitzt jede Maschine $i$ eine
Geschwindigkeit $s_i$, sodass sich die Bearbeitungsdauer als $p_{ij} = p_j /
s_i$ ergibt; die Maschinen unterscheiden sich also nur um einen einzigen Faktor,
und wer bei einem Job schneller ist, ist es bei allen. Bei \Index{unabhängigen
parallelen Maschinen} ($R$) schließlich ist $p_{ij}$ eine beliebige Matrix: Ein
Mitarbeiter kann bei einem Arbeitsgang der schnellste und bei einem anderen der
langsamste sein.

\begin{figure}[tb]
  \centering
  % Einordnung des Problems in die Maschinenmodelle des Scheduling.
% Farben aus nxtlatexcolors.sty -- dieselbe Quelle wie Python und MetaPost.
\begin{tikzpicture}[
    alt={Drei Kästen nebeneinander, von links nach rechts immer allgemeiner:
      identische Maschinen, bei denen alle Balken gleich lang sind;
      uniforme Maschinen, bei denen die Balken je Person um einen Faktor
      verschieden sind; unabhängige Maschinen mit beliebig verschiedenen
      Balken, hervorgehoben als der Fall dieser Arbeit.},
    font=\footnotesize,
    modell/.style={
      draw=nxt.gray2, line width=.6pt, rounded corners=2pt,
      minimum width=32mm, minimum height=20mm, align=center,
      fill=nxt.verylightgray,
    },
    aktiv/.style={modell, draw=nxtgreen, line width=1pt, fill=nxtgreen!8},
    balken/.style={draw=none, fill=nxt.gray3},
    balkenaktiv/.style={draw=none, fill=nxtgreen},
    beschriftung/.style={font=\scriptsize\itshape, text=nxt.gray5},
  ]

  % --- Identische Maschinen: p_ij = p_j -------------------------------
  \node[modell] (id) at (0,0) {};
  \node[above=1mm of id, font=\footnotesize\bfseries] {Identisch};
  \node[below=1mm of id, beschriftung] {$P\,\|\,C_{\max}$};
  \foreach \y/\w in {0/1.0, 1/1.0, 2/1.0} {
    \fill[balken] ($(id.west)+(3mm,-5mm+\y*5mm)$) rectangle ++(\w*22mm, 3mm);
  }
  \node[below=5mm of id, beschriftung, text width=32mm] {$p_{ij}=p_j$\\jeder gleich schnell};

  % --- Uniforme Maschinen: p_ij = p_j / s_i ---------------------------
  \node[modell] (uni) at (4.4,0) {};
  \node[above=1mm of uni, font=\footnotesize\bfseries] {Uniform};
  \node[below=1mm of uni, beschriftung] {$Q\,\|\,C_{\max}$};
  \foreach \y/\w in {0/1.0, 1/0.75, 2/0.5} {
    \fill[balken] ($(uni.west)+(3mm,-5mm+\y*5mm)$) rectangle ++(\w*22mm, 3mm);
  }
  \node[below=5mm of uni, beschriftung, text width=32mm]
    {$p_{ij}=p_j/s_i$\\ein Tempo je Person};

  % --- Unabhängige Maschinen mit Zulässigkeit: unser Fall -------------
  \node[aktiv] (unrel) at (8.8,0) {};
  \node[above=1mm of unrel, font=\footnotesize\bfseries, text=nxtgreen]
    {Unabhängig};
  \node[below=1mm of unrel, beschriftung, text=nxtgreen]
    {$R\,|\,M_j\,|\,C_{\max}$};
  \foreach \y/\w in {0/0.9, 1/0.45, 2/0.7} {
    \fill[balkenaktiv] ($(unrel.west)+(3mm,-5mm+\y*5mm)$) rectangle ++(\w*22mm, 3mm);
  }
  % Gestrichelter Balken = Qualifikation fehlt, Zuordnung unzulässig.
  \draw[draw=nxtgreen, densely dashed, line width=.5pt]
    ($(unrel.west)+(3mm,5mm)$) rectangle ++(6mm, 3mm);
  \node[below=5mm of unrel, beschriftung, text width=34mm, text=nxtgreen]
    {$p_{ij}$ frei, $M_j$ schränkt ein\\ \textbf{der Fall dieser Arbeit}};

  \draw[->, draw=nxt.gray3, line width=.6pt]
    ($(id.east)+(1mm,0)$) -- ($(uni.west)-(1mm,0)$);
  \draw[->, draw=nxt.gray3, line width=.6pt]
    ($(uni.east)+(1mm,0)$) -- ($(unrel.west)-(1mm,0)$);
  \node[beschriftung] at (2.2,1.55) {allgemeiner};
  \node[beschriftung] at (6.6,1.55) {allgemeiner};
\end{tikzpicture}

  \caption[Maschinenmodelle im Vergleich]{Einordnung des Problems: identische, uniforme und unabhängige parallele Maschinen}
  \label{abb:maschinenmodell}
\end{figure}

\subsection{Das Problem \texorpdfstring{$R\mid M_j\mid C_{\max}$}{R|Mj|Cmax}}

Das Problem dieser Arbeit trägt die Signatur $R \mid M_j \mid C_{\max}$. Die
drei Felder sind wie folgt zu lesen.

Das $R$ steht für unabhängige parallele Maschinen. Diese Wahl ist keine
vorsorgliche Verallgemeinerung, sondern durch den Gegenstand erzwungen: Die
betrachteten Mitarbeiter unterscheiden sich nicht in einer pauschalen
Arbeitsgeschwindigkeit, sondern in ihrer Erfahrung mit jeweils bestimmten
Arbeitsgängen. Ein uniformes Modell würde die Daten sichtbar verfälschen.

Das $M_j$ bezeichnet die \Index{Maschinenzulässigkeit}: Für jeden Arbeitsgang
$j$ existiert eine Menge $M_j$ derjenigen Mitarbeiter, die ihn überhaupt
ausführen dürfen. Diese Restriktion bildet die \Index{Qualifikation} ab. Wer für
einen Arbeitsgang nicht qualifiziert ist, kommt für ihn nicht in Betracht --
nicht zu höheren Kosten, sondern gar nicht. Modelltechnisch ist es daher
sauberer und zugleich effizienter, unzulässige Paare erst gar nicht anzulegen,
statt sie über eine große Strafkonstante auszuschließen.

Das $C_{\max}$ schließlich ist der \Index{Makespan}, also der Zeitpunkt, zu dem
auch der letzte Mitarbeiter fertig ist. Bezeichnet $x_{ij} \in \{0,1\}$ die
Entscheidung, ob Mitarbeiter $i$ den Arbeitsgang $j$ übernimmt, so ist
\begin{equation}
  C_{\max} = \max_{i} \; \sum_{j} p_{ij}\, x_{ij}.
  \label{eq:makespan}
\end{equation}
Der Makespan misst nicht die geleistete Arbeit, sondern die verstrichene Zeit.
Er wird allein vom am stärksten belasteten Mitarbeiter bestimmt -- dem
\Index{Engpass}. Diese Eigenschaft ist charakteristisch und hat weitreichende
Folgen: Sie macht den Makespan zu einem sinnvollen Maß für die Frage
\enquote{Wann ist der Auftrag abgeschlossen?}, sie führt aber, wie in
Abschnitt~\ref{sec:sa} zu zeigen sein wird, zu erheblichen Schwierigkeiten bei
Verfahren, die sich von lokalen Verbesserungen leiten lassen.

\section{Komplexität}
\label{sec:komplexitaet}

\subsection{NP-Schwere}

Die naheliegende Idee, alle Zuordnungen aufzuzählen, scheitert an ihrer Anzahl.
Bei $n$ Arbeitsgängen und $m$ zulässigen Mitarbeitern je Arbeitsgang existieren
bis zu $m^n$ Zuordnungen; der Suchraum wächst exponentiell. Schon bei moderaten
Instanzgrößen übersteigt die vollständige Enumeration jede vertretbare
Rechenzeit.

Dass es sich dabei nicht um einen Mangel an Einfallsreichtum handelt, sondern um
eine Eigenschaft des Problems, zeigt die Komplexitätstheorie. Bereits der
Spezialfall $P2 \mid\mid C_{\max}$ -- zwei identische Maschinen, keine
Nebenbedingungen -- ist \Index{NP-schwer}; er enthält das Problem
\Index{Partition}, eines der klassischen NP-vollständigen Probleme im Katalog
von \textcite{garey1979computers}, als Entscheidungsvariante. Da $R \mid M_j \mid
C_{\max}$ diesen Spezialfall umfasst, ist es mindestens ebenso schwer. Nach
gegenwärtigem Kenntnisstand existiert somit kein Algorithmus, der für jede
Instanz in polynomieller Zeit eine beweisbar optimale Lösung liefert.

Es lohnt, die Tragweite dieser Aussage nicht zu überschätzen. NP-Schwere ist
eine Aussage über das Verhalten im schlechtesten Fall und über beliebig große
Instanzen. Sie schließt nicht aus, dass konkrete Instanzen praktischer Größe
exakt und schnell lösbar sind -- moderne Löser profitieren erheblich von
Struktur, insbesondere von der durch $M_j$ erzeugten Ausdünnung des Modells. Die
NP-Schwere ist also kein Verbot, exakte Verfahren einzusetzen; sie ist eine
Warnung davor, sich auf ihre Skalierbarkeit zu verlassen.

\subsection{Approximierbarkeit}

Wenn Optimalität nicht garantiert werden kann, stellt sich die Frage nach
garantierter Nähe zum Optimum. Ein Verfahren heißt
\index{Approximationsalgorithmus}\emph{$\rho$-Approximation}, wenn es in
polynomieller Zeit eine Lösung liefert, deren Zielwert höchstens um den Faktor
$\rho$ über dem Optimum liegt.

Das grundlegende Resultat für unabhängige parallele Maschinen stammt von
\textcite{lenstra1990approximation}. Sie zeigen, dass sich für $R \mid\mid
C_{\max}$ eine \Index{2-Approximation} in polynomieller Zeit konstruieren lässt.
Das Verfahren beruht auf einem \Index{Rundungsverfahren}: Zunächst wird die
\Index{LP-Relaxation} gelöst, deren Lösung Arbeitsgänge anteilig auf mehrere
Maschinen verteilen darf; anschließend werden die gebrochenen Anteile so
gerundet, dass der Fehler beschränkt bleibt. Dieselbe Arbeit zeigt zudem, dass
keine $\rho$-Approximation mit $\rho < 3/2$ existieren kann, sofern nicht
$\mathrm{P} = \mathrm{NP}$ gilt. Zwischen $3/2$ und $2$ klafft bis heute eine
Lücke.

Für die vorliegende Arbeit sind diese Ergebnisse vor allem als Maßstab von
Bedeutung. Sie belegen, dass die Klasse $R$ nicht beliebig unzugänglich ist,
sie zeigen aber auch, dass die theoretischen Garantien großzügig ausfallen: Ein
Faktor $2$ bedeutet im schlechtesten Fall die doppelte Gesamtdauer. Die in
dieser Arbeit eingesetzten Heuristiken tragen, wie sich zeigen wird, überhaupt
keine Gütegarantie -- liefern aber praktisch weit bessere Ergebnisse, als es die
Theorie im schlechtesten Fall zusichern würde.

\section{Exakte Verfahren}
\label{sec:exakt}

\subsection{Ganzzahlige lineare Optimierung}

Ein Modell der \index{Ganzzahlige lineare Optimierung}\emph{ganzzahligen linearen
Optimierung} (\ac{milp}) beschreibt ein Problem durch eine lineare Zielfunktion
und lineare Nebenbedingungen über teils ganzzahligen Variablen
\parencite{wolsey2020integer}. Das hier betrachtete Problem lässt sich so
formulieren:
\begin{align}
  \min \quad & C_{\max} \label{eq:milp-ziel}\\
  \text{u.\,d.\,N.} \quad & \sum_{i \in M_j} x_{ij} = 1 && \text{für alle } j \label{eq:milp-vergabe}\\
  & \sum_{j : i \in M_j} p_{ij}\, x_{ij} \le C_{\max} && \text{für alle } i \label{eq:milp-auslastung}\\
  & x_{ij} \in \{0,1\}. \label{eq:milp-binaer}
\end{align}
Bedingung~\eqref{eq:milp-vergabe} stellt sicher, dass jeder Arbeitsgang genau
einmal und nur an einen qualifizierten Mitarbeiter vergeben wird. Die
Ungleichungen~\eqref{eq:milp-auslastung} verdienen besondere Beachtung: Sie
fordern lediglich, dass $C_{\max}$ mindestens so groß ist wie die Auslastung
jedes einzelnen Mitarbeiters. Erst im Zusammenspiel mit der Minimierung
in~\eqref{eq:milp-ziel} erzwingen sie die Gleichheit aus~\eqref{eq:makespan} --
ein Standardkniff, mit dem sich das nichtlineare Maximum linear ausdrücken
lässt.

\subsection{LP-Relaxation und Schranken}

Lässt man die Ganzzahligkeitsforderung~\eqref{eq:milp-binaer} fallen und
ersetzt sie durch $0 \le x_{ij} \le 1$, so entsteht die \Index{LP-Relaxation}.
Sie ist in polynomieller Zeit lösbar. Da jede zulässige Lösung des
ganzzahligen Problems auch für die Relaxation zulässig ist, liefert deren
Optimalwert eine \index{Schranke!untere}\emph{untere Schranke} für das Optimum. Jede
konkrete Zuordnung wiederum liefert eine \index{Schranke!obere}\emph{obere Schranke}.
Die Differenz zwischen beiden heißt \Index{Dualitätslücke}; solange sie
positiv ist, kann das Optimum irgendwo dazwischen liegen. Schließt sie sich, ist
die Optimalität bewiesen.

Neben der LP-Relaxation existieren einfachere, kombinatorisch begründete
Schranken. Für den Makespan ist etwa
\begin{equation}
  C_{\max} \;\ge\; \frac{1}{m} \sum_{j} \min_{i \in M_j} p_{ij}
  \label{eq:schranke-allgemein}
\end{equation}
gültig: Selbst wenn jeder Arbeitsgang von dem für ihn schnellsten Mitarbeiter
erledigt und die Last vollkommen gleichmäßig verteilt würde, wäre diese Zeit
nicht zu unterbieten. Eine solche Schranke ist in aller Regel nicht erreichbar
-- sie ignoriert, dass die schnellsten Mitarbeiter einander im Weg stehen --,
aber sie ist ohne nennenswerten Aufwand zu berechnen und erlaubt eine erste
Einschätzung der Güte jeder gefundenen Lösung.

\subsection{Branch-and-Bound}

Das Standardverfahren zur exakten Lösung ganzzahliger Modelle ist
\Index{Branch-and-Bound}, das auf \textcite{land1960automatic} zurückgeht und
bei \textcite{nemhauser1988integer} ausführlich dargestellt ist. Es kombiniert
zwei Ideen. Beim \Index{Verzweigen} wird das Problem an einer Variablen mit
gebrochenem Wert in Teilprobleme zerlegt, etwa in die Fälle $x_{ij} = 0$ und
$x_{ij} = 1$; die Vereinigung der Teilprobleme deckt den ursprünglichen
Lösungsraum vollständig ab. Beim \Index{Beschränken} wird für jedes Teilproblem
die Relaxation gelöst. Liegt deren Wert bereits über der besten bisher
bekannten Lösung, so kann das Teilproblem samt allen seinen Nachfolgern
verworfen werden, ohne es zu durchsuchen -- es kann dort nichts Besseres mehr
geben.

Genau in diesem Verwerfen liegt die Kraft des Verfahrens, und genau darin liegt
auch die Bedeutung des Begriffs \enquote{beweisbar optimal}. Terminiert
Branch-and-Bound regulär, so hat es nicht etwa alle Lösungen betrachtet;
vielmehr hat es für jede nicht betrachtete Region gezeigt, dass sie keine
bessere Lösung enthalten \emph{kann}. Das Ergebnis ist damit nicht eine Lösung,
die man nach längerem Suchen nicht mehr verbessern konnte, sondern eine Lösung
mit mathematischem Beweis ihrer Optimalität. Dieser Unterschied ist der
entscheidende Vorzug exakter Verfahren gegenüber allen Heuristiken: Sie liefern
nicht nur eine Antwort, sondern zugleich die Gewissheit, dass es keine bessere
gibt. Erkauft wird er mit einer Laufzeit, die im schlechtesten Fall exponentiell
bleibt. In dieser Arbeit wird dafür der freie \ac{milp}-Löser GLPK eingesetzt
\parencite{glpk}.

\section{Konstruktionsheuristiken}
\label{sec:heuristiken}

Eine \Index{Konstruktionsheuristik} baut eine Lösung in einem Durchgang auf,
ohne sie nachträglich zu verbessern. Sie trifft jede Entscheidung nach einer
festen Regel und nimmt sie nie zurück. Der Reiz solcher Verfahren liegt in ihrer
Geschwindigkeit und Nachvollziehbarkeit; ihr Preis ist das Fehlen jeder
Gütegarantie.

Das \Index{Greedy-Verfahren} betrachtet die Arbeitsgänge in
ihrer Eingangsreihenfolge und weist jeden demjenigen qualifizierten Mitarbeiter
zu, der ihn am frühesten abschließen würde. Es entscheidet also stets lokal
optimal, ohne die Folgen zu bedenken. Der bekannte Schwachpunkt ist die
Behandlung langer Arbeitsgänge: Trifft ein solcher erst spät ein, sind alle
Mitarbeiter bereits belegt, und er verlängert den Makespan unmittelbar.

Die \Index{LPT-Regel} (\ac{lpt}) setzt genau hier an und dreht die Reihenfolge
um: Die Arbeitsgänge werden absteigend nach ihrer Dauer sortiert und in dieser
Reihenfolge vergeben. Die langen Brocken werden zuerst verteilt, solange noch
Spielraum besteht; die kurzen dienen am Ende dem Feinausgleich. Bei unabhängigen
Maschinen ist dabei zu klären, welche Dauer für die Sortierung maßgeblich ist,
da $p_{ij}$ von der Zuordnung abhängt -- in dieser Arbeit wird die kürzestmögliche
Dauer $\min_{i \in M_j} p_{ij}$ verwendet.

Für identische parallele Maschinen hat \textcite{graham1969bounds} für die
LPT-Regel die klassische Schranke
\begin{equation}
  \frac{C_{\max}^{\text{LPT}}}{C_{\max}^{*}} \;\le\; \frac{4}{3} - \frac{1}{3m}
  \label{eq:graham}
\end{equation}
bewiesen. Sie besagt, dass LPT auf $P \mid\mid C_{\max}$ niemals mehr als rund
ein Drittel über dem Optimum liegt -- ein bemerkenswert starkes Ergebnis für ein
derart einfaches Verfahren.

Es ist jedoch entscheidend, festzuhalten, dass diese Schranke im vorliegenden
Fall \emph{nicht} gilt. Grahams Beweis nutzt an zentraler Stelle aus, dass die
Maschinen identisch sind: Nur dann besitzt jeder Job eine wohldefinierte, von
der Maschine unabhängige Dauer $p_j$, nur dann ist die Sortierung nach Dauer
eindeutig, und nur dann lässt sich der Beitrag des zuletzt beendeten Jobs gegen
die Durchschnittslast abschätzen. Bei unabhängigen Maschinen bricht diese
Argumentation zusammen. Die Sortierung nach der kürzestmöglichen Dauer ist eine
Konvention und keine Eigenschaft des Jobs; ein nach ihr kurzer Arbeitsgang kann
für den Mitarbeiter, der ihn tatsächlich erhält, sehr lang sein. Die LPT-Regel
bleibt für $R \mid M_j \mid C_{\max}$ eine plausible und in der Praxis brauchbare
Faustregel -- aber eben nur das. Sie ohne diesen Vorbehalt mit Grahams Schranke
zu versehen, wäre ein Fehlschluss, der in der Anwendungsliteratur nicht selten
anzutreffen ist.

\section{Simulated Annealing}
\label{sec:sa}

\subsection{Herkunft und Grundidee}

Zwischen den schnellen, aber garantielosen Konstruktionsheuristiken und den
exakten, aber teuren Verfahren stehen die \Index{Metaheuristik}en. Sie
verbessern eine vorhandene Lösung schrittweise, ohne den Lösungsraum
systematisch zu erschöpfen. Das in dieser Arbeit eingesetzte Verfahren ist
\Index{Simulated Annealing} (\ac{sa}).

Sein Name und seine Funktionsweise stammen aus der Metallurgie. Beim
\Index{Ausglühen} eines Werkstücks wird das Metall stark erhitzt und
anschließend langsam abgekühlt. Bei hoher Temperatur besitzen die Atome genug
Energie, um ihre Gitterplätze zu verlassen; mit sinkender Temperatur ordnen sie
sich zunehmend in einem energiearmen, regelmäßigen Kristallgitter an. Schreckt
man das Werkstück dagegen rasch ab, gefrieren die Fehlstellen ein, und das
Material bleibt spröde. Die langsame Abkühlung ist es, die den energetisch
günstigen Zustand ermöglicht.

\textcite{metropolis1953equation} entwickelten ein Verfahren, um dieses
thermodynamische Verhalten zu simulieren. Unabhängig voneinander erkannten
\textcite{kirkpatrick1983optimization} und \textcite{cerny1985thermodynamical},
dass sich das Vorgehen auf beliebige Optimierungsprobleme übertragen lässt, wenn
man die Zielfunktion als \Index{Energie} und einen Kontrollparameter als
\Index{Temperatur} auffasst. Eine ausführliche theoretische Behandlung findet
sich bei \textcite{laarhoven1987simulated}.

\subsection{Das Metropolis-Kriterium}

Der Kern des Verfahrens ist die Regel, nach der über die Annahme eines
Kandidaten entschieden wird. Aus der aktuellen Lösung wird eine Nachbarlösung
erzeugt und die Energiedifferenz $\Delta E$ bestimmt. Ist $\Delta E \le 0$, die
Lösung also nicht schlechter, wird sie stets angenommen. Ist sie schlechter,
wird sie nicht etwa verworfen, sondern mit der Wahrscheinlichkeit
\begin{equation}
  P(\text{Annahme}) = \exp\!\left(-\frac{\Delta E}{T}\right)
  \qquad \text{für } \Delta E > 0
  \label{eq:metropolis}
\end{equation}
dennoch übernommen. Dieses \Index{Metropolis-Kriterium} ist der eigentliche
Grund für die Leistungsfähigkeit des Verfahrens. Ein Verfahren, das nur
Verbesserungen akzeptiert, bleibt im ersten \index{Optimum!lokales}\emph{lokalen
Optimum} stehen, das es erreicht. Die gelegentliche Annahme einer
Verschlechterung erlaubt es, ein solches Tal wieder zu verlassen.

Die Temperatur $T$ steuert dabei die Bereitschaft zur Verschlechterung. Bei
großem $T$ geht der Exponent gegen null und die Wahrscheinlichkeit gegen eins:
Das Verfahren nimmt fast alles an und irrt nahezu zufällig durch den Suchraum.
Bei kleinem $T$ geht die Wahrscheinlichkeit gegen null, und das Verfahren
verhält sich wie eine reine lokale Suche. Bemerkenswert ist ferner, dass
\eqref{eq:metropolis} nicht nur von $T$, sondern auch von der Größe von $\Delta
E$ abhängt: Kleine Verschlechterungen werden deutlich bereitwilliger
hingenommen als große.

\subsection{Abkühlung}

Die \Index{Abkühlung} überführt das Verfahren von der Exploration in die
Verfeinerung. Praktisch verbreitet ist die \index{Abkühlung!geometrische}\emph{geometrische Abkühlung}
\begin{equation}
  T_{k+1} = \alpha \cdot T_k, \qquad 0 < \alpha < 1,
  \label{eq:abkuehlung}
\end{equation}
bei der die Temperatur in jedem Schritt um einen festen Faktor sinkt und somit
$T_k = T_0 \cdot \alpha^k$ gilt. Gibt man Anfangstemperatur $T_0$,
Endtemperatur $T_N$ und Schrittzahl $N$ vor, so ergibt sich $\alpha =
(T_N/T_0)^{1/N}$. Diese Parametrisierung ist der direkten Wahl von $\alpha$
vorzuziehen, weil sie an inhaltlich interpretierbaren Größen ansetzt.

Theoretisch lässt sich zeigen, dass \ac{sa} unter einer hinreichend langsamen,
logarithmischen Abkühlung mit Wahrscheinlichkeit eins gegen das globale Optimum
konvergiert \parencite{laarhoven1987simulated}. Dieses Resultat ist praktisch
allerdings ohne Wert, da die erforderlichen Laufzeiten die vollständige
Enumeration übersteigen. In der Anwendung kühlt man schneller ab und verzichtet
damit auf jede Garantie -- \ac{sa} liefert eine gute Lösung, aber keinen Beweis.

\subsection{Nachbarschaft}

Die \Index{Nachbarschaft} legt fest, welche Lösungen aus einer gegebenen Lösung
in einem Schritt erreichbar sind. Ihre Gestaltung ist die eigentliche
problemspezifische Entwurfsentscheidung: Ist sie zu klein, zerfällt der
Suchraum in unerreichbare Gebiete; ist sie zu groß, verliert der Schritt seinen
lokalen Charakter und das Verfahren entartet zur Zufallssuche.

In dieser Arbeit werden die beiden in Abbildung~\ref{abb:nachbarschaft}
dargestellten Operatoren kombiniert. Das \index{Nachbarschaft!Verschieben}\emph{Verschieben}
nimmt einem Mitarbeiter einen Arbeitsgang ab und übergibt ihn einem anderen
qualifizierten Mitarbeiter. Das \index{Nachbarschaft!Tauschen}\emph{Tauschen} kreuzt
zwei Arbeitsgänge zwischen ihren Mitarbeitern, sofern beide für den jeweils
anderen Arbeitsgang qualifiziert sind. Beide werden benötigt: Das Verschieben
allein genügt nicht, denn sind zwei Mitarbeiter gleichermaßen ausgelastet, so
verlagert jedes Verschieben das Problem nur, während allein ein Tausch die Last
umverteilen kann, ohne sie an anderer Stelle zu erhöhen.

\begin{figure}[tb]
  \centering
  % Die beiden Nachbarschaftsoperatoren des Verfahrens.
% Farben aus nxtlatexcolors.sty.
\begin{tikzpicture}[
    alt={Zwei Skizzen. Links Verschieben: Eine Aufgabe des Engpasses
      wandert von Dobrev zu Haas. Rechts Tauschen: Bauer und Ivanov tauschen
      je eine Aufgabe, was nur geht, wenn beide für die Aufgabe des anderen
      qualifiziert sind.},
    font=\footnotesize,
    person/.style={font=\scriptsize, text=nxt.gray5, anchor=east},
    aufgabe/.style={
      draw=none, rounded corners=1pt, minimum height=4mm,
      font=\scriptsize\color{white}, inner sep=1pt, anchor=west,
    },
    pfeil/.style={->, draw=nxt.gray4, line width=.7pt, shorten >=1pt},
    titel/.style={font=\footnotesize\bfseries},
    hinweis/.style={font=\scriptsize\itshape, text=nxt.gray5},
  ]

  % ============ links: Verschieben =====================================
  \begin{scope}
    \node[titel] at (1.9,1.5) {Verschieben};
    \node[hinweis, text width=38mm, align=center] at (1.9,-1.35)
      {eine Aufgabe des Engpasses\\an eine andere Person};

    \node[person] at (0,0.8) {Dobrev};
    \node[person] at (0,0.0) {Haas};

    % Dobrev ist der Engpass: langer Balken.
    \node[aufgabe, fill=nxt.gray3, minimum width=11mm] at (0.15,0.8) {C-310};
    \node[aufgabe, fill=nxtgreen, minimum width=9mm]  at (1.35,0.8) {C-333};
    \node[aufgabe, fill=nxt.gray3, minimum width=8mm] at (2.45,0.8) {D-042};
    \node[aufgabe, fill=nxt.gray3, minimum width=7mm] at (0.15,0.0) {F-078};

    \draw[pfeil] (1.8,0.6) to[out=-90,in=90] (1.35,0.2);
    \node[aufgabe, fill=nxtgreen, minimum width=9mm, opacity=.28]
      at (1.05,0.0) {};
  \end{scope}

  % ============ rechts: Tauschen =======================================
  \begin{scope}[xshift=62mm]
    \node[titel] at (1.9,1.5) {Tauschen};
    \node[hinweis, text width=40mm, align=center] at (1.9,-1.35)
      {zwei Aufgaben kreuzen --\\nur wenn beide qualifiziert sind};

    \node[person] at (0,0.8) {Bauer};
    \node[person] at (0,0.0) {Ivanov};

    \node[aufgabe, fill=nxtgreen, minimum width=12mm] at (0.15,0.8) {A-115};
    \node[aufgabe, fill=nxt.gray3, minimum width=7mm] at (1.45,0.8) {A-134};
    \node[aufgabe, fill=nxt.orange, minimum width=8mm] at (0.15,0.0) {A-127};
    \node[aufgabe, fill=nxt.gray3, minimum width=9mm] at (1.15,0.0) {B-245};

    \draw[pfeil] (0.75,0.6) to[out=-70,in=70] (0.55,0.2);
    \draw[pfeil] (0.55,0.2) to[out=110,in=-110] (0.75,0.6);
  \end{scope}

  % Trennlinie
  \draw[draw=nxt.gray1, line width=.5pt] (5.0,-1.6) -- (5.0,1.6);
\end{tikzpicture}

  \caption[Nachbarschaftsoperatoren]{Die beiden Nachbarschaftsoperatoren: Verschieben und Tauschen}
  \label{abb:nachbarschaft}
\end{figure}

\subsection{Das Plateauproblem}

Abschließend ist auf eine Schwierigkeit hinzuweisen, die sich unmittelbar aus
der Struktur des Makespans ergibt und für die Auslegung des Verfahrens von
erheblicher Bedeutung ist. Nach \eqref{eq:makespan} ist $C_{\max}$ ein Maximum.
Er ändert sich folglich nur dann, wenn ein Zug gerade den Engpass-Mitarbeiter
betrifft. Jede Umverteilung zwischen zwei nicht ausgelasteten Mitarbeitern
lässt die Zielfunktion vollkommen unverändert.

Verwendet man den Makespan unmittelbar als Energie, so ist die Suchlandschaft
daher von ausgedehnten \index{Plateau}\emph{Plateaus} durchzogen: Weite Bereiche
besitzen exakt denselben Zielwert, und $\Delta E = 0$ liefert dem Verfahren
keinerlei Hinweis, welche Richtung lohnend ist. Das Verfahren irrt über diese
Ebenen, ohne geführt zu werden. Dieses \Index{Plateauproblem} tritt bei allen
Max- und Min-Zielfunktionen auf und ist keine Eigenheit des vorliegenden Falls.

Ein bewährter Ausweg besteht darin, die Energie um einen kleinen Zusatzterm zu
ergänzen, der auch dann reagiert, wenn der Makespan es nicht tut. Bezeichnet
$a_i$ die Auslastung des Mitarbeiters $i$, so bietet sich
\begin{equation}
  E = \max_i a_i \;+\; \lambda \sqrt{\sum_i a_i^{\,2}}
  \label{eq:glaettung}
\end{equation}
an. Der zweite Term, die euklidische Norm des Auslastungsvektors, sinkt bereits
dann, wenn die Last gleichmäßiger verteilt wird -- auch ohne dass der Makespan
fällt. Damit erhält das Verfahren auf dem Plateau ein Gefälle, dem es folgen
kann. Entscheidend ist die Wahl des Gewichts $\lambda$: Es muss klein genug
sein, dass der Zusatzterm die Rangfolge zweier Lösungen mit verschiedenem
Makespan niemals umkehrt. Er dient allein der Wegfindung, nicht der Bewertung.
Die eigentliche Zielgröße bleibt der Makespan, und ausschließlich nach ihm wird
die beste gefundene Lösung ausgewählt und berichtet. Welchen Unterschied diese
Modifikation praktisch ausmacht, wird im Ergebniskapitel quantifiziert.
